\section{Gaussian tubes and the least singular value}
\label{sec:gaussian-tubes}

Proposition~\ref{prop:distance} reduces the small singular value event to proximity to a real polynomial zero set.
We first bound the Gaussian measure of the points within a given distance of such a set.

\begin{lemma}
\label{lem:gaussian-tube}
Let \(q\) be a nonzero, nonnegative real polynomial on \(\mathbb R^n\) of degree at most \(D\), where \(D\ge 1\), and suppose \(Z:=q^{-1}(0)\) is nonempty.
If \(g\) is a standard Gaussian vector in \(\mathbb R^n\), then for every \(\rho>0\),
\[
  \mathbb P\bigl\{\operatorname{dist}(g,Z)\le \rho\bigr\}
  \le C D\sqrt n\,\rho\,
      \exp\bigl(C(D+1)\sqrt n\,\rho\bigr)+e^{-cn}
\]
for universal constants \(C,c>0\).
\end{lemma}

\begin{proof}
Let \(X\) be uniform in the ball of radius \(\sigma>0\) centered at the origin in \(\mathbb R^n\).
For a smooth hypersurface \(Y\) given by a regular level of a polynomial of degree at most \(D\), the codimension-one case of Lotz's tube estimate \cite[Theorem~1.1]{Lot15} gives
\[
  \mathbb P\bigl\{\operatorname{dist}(X,Y)\le\rho\bigr\}
  \le 4\sum_{r=1}^n \binom nr
       (2D\rho/\sigma)^r(1+\rho/\sigma)^{n-r}.
\]
Set \(u:=\rho/\sigma\), so the right side is \(4((1+(2D+1)u)^n-(1+u)^n)\).
By the mean value theorem, this expression is bounded by
\begin{equation}
  \frac{8nD\rho}{\sigma}\exp\!\left(\frac{n(2D+1)\rho}{\sigma}\right).
  \label{eq:uniform-tube}
\end{equation}

To apply this bound to \(Z\), fix \(\eta>0\) and consider the zeros of \(q\) in the closed ball of radius \(\sigma+\rho\).
Since \(Z\) is closed, every point of the tube in the ball of radius \(\sigma\) has a nearest zero in this larger ball.
If there are no such zeros, the tube event in the ball of radius \(\sigma\) is empty.
Otherwise, near each such zero \(y\) there is a point \(p_y\) with \(|p_y-y|<\eta/2\) and \(q(p_y)>0\), since a nonzero nonnegative polynomial cannot vanish on an open ball.
By compactness, finitely many open balls of radius \(\eta/2\) centered at these zeros cover the relevant zero set.
Choose a positive regular value \(\delta\) of \(q\) less than all the values \(q(p_y)\) for the chosen centers, using Sard's theorem.
For each relevant zero \(y'\), the segment from \(y'\) to the positive point associated with a covering ball lies within distance \(\eta\) of \(y'\) and meets \(q^{-1}(\delta)\).
Every point in the ball of radius \(\sigma\) within distance \(\rho\) of \(Z\) is therefore within distance \(\rho+\eta\) of \(q^{-1}(\delta)\).
The latter is a smooth hypersurface defined by a polynomial of degree at most \(D\).
Applying \eqref{eq:uniform-tube} with \(\rho+\eta\) and letting \(\eta\) decrease to zero gives \eqref{eq:uniform-tube} for \(Z\).

For the Gaussian estimate, first sample \(S>0\) with the chi distribution with \(n+2\) degrees of freedom and then sample a point uniformly in the ball of radius \(S\).
With \(\omega_n\) the volume of the unit ball in \(\mathbb R^n\), the resulting density at \(x\in\mathbb R^n\) is
\[
  \int_{|x|}^\infty
    \frac{1}{\omega_n s^n}
    \frac{s^{n+1}e^{-s^2/2}}{2^{n/2}\Gamma(n/2+1)}\,ds
  = (2\pi)^{-n/2}e^{-|x|^2/2}.
\]
Since \(S^2\) is the sum of \(n+2\) independent squared standard normal variables, the exponential Markov bound gives
\(\mathbb P\{S<\sqrt n/2\}\le e^{n/8}2^{-(n+2)/2}\le e^{-cn}\) for a universal \(c>0\).
On \(S\ge\sqrt n/2\), the bound \eqref{eq:uniform-tube} for \(Z\) has prefactor at most \(16D\sqrt n\,\rho\) and exponent at most \(2(2D+1)\sqrt n\,\rho\).
Together these estimates give the claim.
\end{proof}

\begin{proof}[Proof of Theorem~\ref{thm:pencil}]
Write each component of \(H^{(N)}\) in a real Hilbert--Schmidt orthonormal basis of Hermitian matrices.
Its law in these coordinates is that of \(N^{-1/2}g\), where \(g\) is standard Gaussian in \(\mathbb R^n\) and \(n:=dN^2\).
Indeed, \(\Tr(H_j^2)\) is the squared Euclidean norm of the coordinate vector.
Using the same coordinates to evaluate the pencil, define
\[
  q_N(x):=\bigl|\det T_N(N^{-1/2}x)\bigr|^2,
  \qquad Z_N:=q_N^{-1}(0).
\]
The determinant has degree at most \(mN\), so \(q_N\) is a nonnegative real polynomial of degree at most \(2mN\), with \(q_N(0)=1\).

Fix \(M>0\) and let \(C_0,\alpha,\varepsilon_0>0\) be the constants of Proposition~\ref{prop:distance}.
If \(Z_N\) is empty, the proposition excludes \(s_{\min}(T_N(H^{(N)}))\le\varepsilon\) on \(\max_j\|H_j^{(N)}\|\le M\) for \(0<\varepsilon\le\varepsilon_0\).
Otherwise, on the same event the proposition gives
\[
  s_{\min}(T_N(H^{(N)}))\le\varepsilon
  \quad\Longrightarrow\quad
  \operatorname{dist}(g,Z_N)\le C_0\sqrt N\,\varepsilon^\alpha.
\]
The factor \(\sqrt N\) converts the Hilbert--Schmidt distance in the original matrix coordinates to distance in the standard Gaussian coordinates.
Apply Lemma~\ref{lem:gaussian-tube} with \(D:=2mN\) and \(\rho:=C_0\sqrt N\,\varepsilon^\alpha\).
Since \(\sqrt n=\sqrt d\,N\), the resulting probability bound is
\[
  \mathbb P\bigl\{s_{\min}(T_N(H^{(N)}))\le\varepsilon,\
                 \max_j\|H_j^{(N)}\|\le M\bigr\}
  \le C N^{5/2}\varepsilon^\alpha
      \exp\bigl(CN^{5/2}\varepsilon^\alpha\bigr)+e^{-cN^2}.
\]
Take \(M>2\), so that the norm event has probability tending to one, and choose \(A>0\) with \(A\alpha>5/2\).
The right side tends to zero for \(\varepsilon=N^{-A}\), proving \eqref{eq:pencil-floor}.
\end{proof}
